% =============================================================================
% AGENTE  : Cientista-Líder
% MÓDULO  : Conclusão — Português Brasileiro (conclusion_pt.tex)
% SINCRON.: Deve ser consistente com conclusion_en.tex
% STATUS  : RASCUNHO
% =============================================================================

\section{Conclusão}
\label{sec:conclusion}

Este estudo demonstra um pipeline reprodutível e multi-agente para a
caracterização estatística e o relato científico de painéis de marcadores
moleculares. O conjunto de dados simulado confirmou que os marcadores \ssr{}
apresentam valores de \pic{} e heterozigose esperada superiores por locus
em comparação aos marcadores \snp{}, enquanto estes últimos oferecem
vantagens em termos de densidade de marcadores e escalabilidade da
genotipagem.

A estrutura modular do documento \LaTeX{}, combinada com scripts de análise
Python automatizados, permite iteração rápida à medida que novos dados se
tornem disponíveis ou que os requisitos de formatação do journal mudem.
Este arcabouço será aplicado a dados de genotipagem empírica de [espécie
alvo] nas fases subsequentes deste programa de pesquisa.
